\documentclass[10pt]{amsart}
\usepackage{amsmath,amssymb,amsthm}
% US letter, margins of about 1.8 cm set by hand; compact section headings.
\setlength{\textwidth}{\paperwidth}\addtolength{\textwidth}{-3.6cm}
\setlength{\oddsidemargin}{-0.74cm}\setlength{\evensidemargin}{-0.74cm}
\setlength{\headheight}{0pt}\setlength{\headsep}{0pt}
\setlength{\topmargin}{-0.84cm}
\setlength{\textheight}{\paperheight}\addtolength{\textheight}{-3.6cm}
\setlength{\footskip}{0.8cm}
\ifdefined\pdfpagewidth\pdfpagewidth=\paperwidth\pdfpageheight=\paperheight\fi
\pagestyle{plain}\raggedbottom % no stretched pages, hence no underfull pages
\setlength{\parskip}{0pt plus 0.5pt}
\widowpenalty=150 \clubpenalty=150
\makeatletter\def\section{\@startsection{section}{1}{\z@}{.6\baselineskip}{.25\baselineskip}{\normalfont\bfseries}}\makeatother

\numberwithin{equation}{section}
\newtheoremstyle{compact}{4pt}{4pt}{\itshape}{}{\bfseries}{.}{.5em}{}
\theoremstyle{compact}
\newtheorem{theorem}{Theorem}[section]
\newtheorem{proposition}[theorem]{Proposition}
\newtheorem{lemma}[theorem]{Lemma}
\newtheorem{corollary}[theorem]{Corollary}
\newtheorem{conjecture}[theorem]{Conjecture}
\newtheorem*{mainthm}{Main Theorem}
\makeatletter % the proof environment of amsthm, with the spacing of the theorems
\renewenvironment{proof}[1][\proofname]{\par\pushQED{\qed}\normalfont
  \topsep4\p@\relax\trivlist\item[\hskip\labelsep\itshape#1\@addpunct{.}]%
  \ignorespaces}{\popQED\endtrivlist\@endpefalse}
\makeatother

\newcommand{\CP}{\mathbb{CP}}
\newcommand{\CPb}{\overline{\mathbb{CP}}{}^2}
\newcommand{\Z}{\mathbb Z}
\newcommand{\Q}{\mathbb Q}
\newcommand{\F}{\mathbb F}
\newcommand{\PD}{\operatorname{PD}}
\newcommand{\Sym}{\operatorname{Sym}}
\newcommand{\SO}{\mathrm{SO}}
\newcommand{\fs}{\mathfrak s}
\newcommand{\ft}{\mathfrak t}
\newcommand{\cO}{\mathcal O}
\newcommand{\cV}{\mathcal V}
\newcommand{\cM}{\mathcal M}
\newcommand{\cZ}{\mathcal Z}
\newcommand{\hS}{\hat\Sigma}
\newcommand{\step}[1]{\par\noindent\textit{#1}\enspace}
\newcommand{\hyp}[1]{\par\noindent\textup{#1}\enspace}
% labels of claims
\newcommand{\ver}{\textsc{verified}}
\newcommand{\pla}{\textsc{plausible}}
\newcommand{\spe}{\textsc{speculative}}
% References set as one paragraph, to keep the note within five pages.
\makeatletter
\renewenvironment{thebibliography}[1]{\par\footnotesize\noindent
  \textbf{References.}\ignorespaces
  \def\@lbibitem[##1]##2{\unskip\hskip.5em plus.3em minus.1em[##1]~\if@filesw\immediate
    \write\@auxout{\string\bibcite{##2}{##1}}\fi\ignorespaces}%
  \def\bysame{---}}{\par}
\makeatother

\begin{document}

{\centering\bfseries FOUR STATEMENTS\par}

\section*{Overview}

Let $K\subset S^3$ be a knot, $Y_r=S^3_r(K)$, and suppose that
$\phi\colon Y_{-2}\to Y_2$ is an orientation-preserving diffeomorphism. The
Strategy computes one integer in two ways: a count $\Omega$ of instantons over
a cube of metrics on a closed four-manifold $X(\Pi)$, assembled from two caps
and the cobordisms of a composite $\Pi$ of the maps $\phi_*B$ and $HB$. The four
statements below are its four steps. In \S1 (Steps~2 and~4 of the Strategy)
the map $B$ of the negative-definite cobordism $W_B\colon Y_2\to Y_{-2}$, defined
by an interval of metrics and cut down by a sphere of square $-4$, is an
integral chain map; a relation for $(-4)$-spheres and a constraint on the
gradings of a cosmetic pair come with it. In \S2 (Step~2), if $B$ is a rational
isomorphism, the Floer pairing of two caps built from the symplectic manifold
of Kronheimer and Mrowka is nonzero along infinitely many composites $\Pi_M$.
In \S3 (Step~4) $\Omega(X(\Pi))$ equals that pairing. In \S4 (Steps~1, 3, 5
and~6) $\Omega(X(\Pi_M))=0$ for large $M$, by Stokes' theorem on a cut-down
moduli space of $\SO(3)$-monopoles over the cube. We assume:

\hyp{(H1)} Floer's instanton homology of $Y_{\pm2}$, whose only reducible flat
connections are the central ones $\theta$, $\chi\theta$, is defined over $\Z$
with the absolute $\Z/8$-grading fixed by $\theta$; a cobordism with $b_1=0$
carrying a class $c$ that vanishes on its ends induces a map of degree
$-2c^2-3b^+\pmod 8$ \cite[proof of Prop.~3.4]{DLME}; Donaldson invariants glue
along $Y_{\pm2}$ and along $(Y_0,w)$ as Floer pairings; and the meridional
handle cobordisms $Y_{-2}\to Y_{-1}$, $Y_1\to Y_2$ induce isomorphisms over
$\F_2$ (those through $Y_0$ do so over $\Z$ by Floer's exact triangle
\cite[Prop.~3.8]{DLME}, as $I(S^3)=0$).
\hyp{(H2)} $B\otimes\Q$ is an isomorphism (Conjecture~\ref{conj:B}).
\hyp{(H3)} Witten's formula holds for the closed symplectic manifold $X_0$ of
\cite[Prop.~15]{KM04}. This is \cite[Cor.~7]{KM04}, whose proof rests on
\cite[Thm.~1.1]{FLM}, stated there under the gluing hypothesis
\cite[Hyp.~7.8.1]{FLM}.
\hyp{(H4)} Gluing in families of broken metrics at irreducible nondegenerate
flat connections, and at the flat connection $\gamma$ of trace zero on
$L(4,1)$ against the reducible of energy $\frac14$ on $\nu S$, with
orientations over $\Z$; in particular the relative sign $\epsilon$ of
Theorem~\ref{thm:B} is constant.
\hyp{(H5)} The analytic foundations of Theorem~\ref{thm:family} listed in \S4,
uniformly in the neck length.

\begin{mainthm}
Let $K\subset S^3$ be a knot of genus at least two, and assume
\textup{(H1)--(H5)}. Then there is no orientation-preserving diffeomorphism
$S^3_{-2}(K)\to S^3_2(K)$.
\end{mainthm}

\begin{proof}
For the composites $\Pi_M$ of Corollary~\ref{cor:composites}, that corollary and
Theorem~\ref{thm:glue} give $\Omega(X(\Pi_M))\ne0$ for infinitely many
$M\in8\Z$, and Theorem~\ref{thm:noreducible} gives $\Omega(X(\Pi_M))=0$ for all
large $M\in8\Z$, for the same manifolds, metrics and divisors.
\end{proof}

By Ni--Wu \cite{NW} and Hanselman \cite[Thm.~2]{Ha}, a cosmetic pair in $S^3$ is
$\pm1/n$ or $\pm2$, and $\pm2$ forces genus two; Daemi, Lidman and Miller
Eismeier exclude $\pm1/n$ \cite[Thm.~1.3, Cor.~1.4]{DLME}. So the Main Theorem
would settle the cosmetic surgery conjecture in $S^3$ under (H1)--(H5)
(\ver\ as logic). Three corrections to the Strategy are built in: rational
invertibility of $B$ replaces the congruence $u^M\equiv1\pmod{2^N}$ of Step~2
(Corollary~\ref{cor:composites}); as $HB$ has degree four
(Proposition~\ref{prop:WH}), what is used, and expected, is only that $B$ is an
isomorphism modulo two; and the Chern--Simons filtration is not known to be
favourable for $B$, rather than known to be unfavourable (\S1).

\section{The distance-four map}

Let $X^\circ_r$ be the trace of $r$-surgery on $K$ with a ball removed, and
$W_B=(-X^\circ_2)\cup_JX^\circ_{-2}\colon Y_2\to Y_{-2}$, $J\cong S^3$ (the $W'$
of the Strategy). It is simply connected with form $\langle-2\rangle^2$ on the
capped Seifert surfaces $F_l$, $F_r$. The cores of the two $2$-handles form a
sphere $S$ with $[S]=F_l-F_r$, $S\cdot S=-4$, $S\cap J=K$ and $\partial\nu
S\cong L(4,1)$ (\ver). Let $G=[-1,1]$ parametrize metrics broken along $J$ at $-1$ and
along $\partial\nu S$ at $+1$ \cite[\S3.9]{KM11}. For a $\mathrm U(2)$-bundle $E$
with $\langle c_1(E),S\rangle$ even, the jumping-line divisor $V_S$ is the set
of connections $A$ for which the holomorphic bundle $(E|_S\otimes(\det
E|_S)^{-1/2},\bar\partial_A)$ is $\cO(k)\oplus\cO(-k)$ with $k\ge1$: the zero
set of the canonical section of the determinant line of $\bar\partial_A$ on its
twist by $\cO(-1)$, of first Chern class $\pm\mu(S)$ \cite{D90}. On $Y_{\pm2}$,
$\chi$ is the flat real line bundle of holonomy $-1$, and $\chi_*$ the
involution it induces on Floer complexes (for $\chi$-invariant perturbations,
\pla); it has degree $4$ and shifts the Chern--Simons functional by $\frac12$, so
it fixes no generator (\ver).

At a fixed metric, cutting the map of $W_B$ down by $V_S$ gives a chain map of
degree $-2$. It is null-homotopic in two ways: by stretching $J$, since the
only flat connection on $S^3$ is trivial, with a three-dimensional stabilizer;
and, for the signed sum over the bundles with $c_1=c$ and $c+\PD(S)$, by
stretching $L(4,1)$, since the two bundles agree off $\nu S$ and the reducible
of energy $\frac14$ on $\nu S$ enters both with opposite signs. The difference
of the two null-homotopies is a chain map:

\begin{theorem}\label{thm:B}
Assume \textup{(H1)} and \textup{(H4)}. For $c\in H^2(W_B;\Z)$ vanishing on
$Y_{\pm2}$, let $m_{G,c}\colon C(Y_2;\Z)\to C(Y_{-2};\Z)$ count the pairs
$(t,[A])$, $t\in\operatorname{int}G$, with $A$ a $g_t$-instanton of index one
and irreducible limits on the bundle $E_c$ with $c_1(E_c)=c$, modulo
determinant-one gauge, and $[A]\in V_S$. There is a sign $\epsilon$, depending
only on orientation conventions, such that $B_c=m_{G,c}-\epsilon\,m_{G,c+\PD(S)}$
is a chain map of degree $-1-2c^2\pmod 8$, whose chain homotopy class does not
depend on the choices. Moreover $B_{c+2a}\simeq\pm\chi_*^{a(F_r)}B_c\,
\chi_*^{a(F_l)}$; so up to sign and homotopy there are exactly two such maps,
$B=B_0$ of degree $-1$ and $B'=B_{-\PD(F_l)}\simeq\pm B\chi_*$ of degree $3$.
\end{theorem}

\begin{proof}[Sketch]
For $\dim G=1$ the boundary formula of \cite[\S3.9]{KM11} reads $\partial
m_{G,c}+m_{G,c}\partial=\pm(m_{1,c}-m_{-1,c})$, where $m_{\pm1,c}$ counts
index-two solutions in $V_S$ for the broken metric at $\pm1$; bubbles and
breaking at $\theta$, $\chi\theta$ are excluded by index. By index additivity
\cite[proof of Prop.~3.4]{DLME}, a solution broken along $J$ with irreducible
halves has index $\operatorname{ind}A_l+3+\operatorname{ind}A_r\ge3$, and a
reducible half has index at least $-3$ \cite[Lemma~5.6]{DLME} and forces a
further central matching on $Y_{\pm2}$; so $m_{-1,c}=0$. On $\nu S$, with
$\xi(S)=1$, the reducibles with $w_2=0$ have $v=2j\xi$, energy $j^2/4$ and
framed index $8E=2j^2$; those with $j=0$ miss $V_S$, as they restrict to
$\cO\oplus\cO$. As the index is $2=\operatorname{ind}(\text{exterior})+8E$, the
solutions at $+1$ are rigid irreducible instantons on $W_B\setminus\nu S$ with
limit $\gamma$, glued to the reducible with $j=\pm1$; it meets $V_S$ once, with
local degree $\pm1$, for its decaying deformations are the extensions of
$\cO(1)$ by $\cO(-1)$, trivial on $S$ unless split. Hence $m_{1,c}=\epsilon_c\,
m_\circ$, where $m_\circ$ counts those instantons and does not depend on $c$, as
$\PD(S)$ vanishes off $\nu S$; and (H4) gives $\epsilon_{c+\PD(S)}=
\epsilon\,\epsilon_c$. Every $c$ is some $2a$, and $E_{c+2a}\cong E_c\otimes
L_a$, $L_a$ flat near the ends with holonomy $\chi^{a(F_l)}$, $\chi^{a(F_r)}$. The
degree is $-2c^2-(\operatorname{codim}V_S-\dim G)$.
\end{proof}

In the inequality of \cite[Lemma~2.7(b)]{DLME}, $B$ has $L-D/8=\frac18-\eta$,
$\eta$ the least energy of a counted instanton \cite[Prop.~3.4]{DLME}; this is
favourable only if $\eta>\frac18$, which is not known (\ver), and is why
monopoles replace the Chern--Simons filtration. The closed counterpart of
Theorem~\ref{thm:B} continues the relations of Kotschick, Ruberman and
Fintushel--Stern for spheres of square $-1$, $-2$, $-3$ \cite[Lemma~2.3(6),
Thms.~2.1, 2.4]{FS1}:

\begin{proposition}\label{prop:minus4}
Let $X$ be closed, oriented and simply connected with $b^+(X)\ge2$, let $S\subset
X$ be an embedded sphere with $S\cdot S=-4$, and $w\in H^2(X;\Z)$ with $w\cdot S$
even. Then $D^{w+\PD(S)}_X(Sz)=-D^w_X(Sz)$ for every $z$ in the algebra
generated by the point class and $S^\perp$.
\end{proposition}

\begin{proof}[Proof for simple type]
Write $\mathbf D^w_X=e^{Q_X/2}\sum_\beta(-1)^{(w^2+\beta\cdot w)/2}a_\beta e^\beta$,
the sum over the basic classes \cite[Thm.~4.1]{FS2}. As $(w+\PD
S)^2=w^2+2w\cdot S-4$ and $w\cdot S$ is even, replacing $w$ by $w+\PD(S)$
multiplies the term of $\beta$ by $(-1)^{\beta\cdot S/2}$, and differentiating
along $S$ at $h\in S^\perp$ multiplies it by $\beta\cdot S$. So
$D^{w+\PD(S)}(S\,\cdot)+D^w(S\,\cdot)$ involves only classes with $\beta\cdot
S\equiv0\pmod4$, $\beta\cdot S\ne0$; by \cite[Thm.~4.3]{FS2} these have
$\beta\cdot S=\pm4$, and $\sum_{\beta\cdot S=4}a_\beta
e^{\beta+S}=(-1)^{(1+b^+)/2}\sum_{\beta\cdot S=4}a_\beta e^{-\beta-S}$. With
$a_{-\beta}=(-1)^{(1+b^+)/2}a_\beta$, which follows from the parity of the
degrees in which $D^w$ is nonzero, this says that the classes with $\beta\cdot
S=-4$ are the $\beta+2\PD(S)$ with $\beta\cdot S=4$, with $a_{\beta+2\PD
S}=a_\beta$. Such a pair has the same sign and the same restriction to
$S^\perp$, and opposite $\beta\cdot S$, so its terms cancel.
\end{proof}

\begin{proposition}\label{prop:WH}
Let $W_H\colon Y_{-2}\to Y_2$ be the composite of the four meridional handle
cobordisms through $Y_{-1}$, $Y_0$, $Y_1$. Every class $c$ on $W_H$ that
vanishes on $Y_{\pm2}$ and is odd on the capped Seifert surface $F_0\subset Y_0$
has $c^2\equiv0$ or $2\pmod4$, so $(W_H,c)_*$ has degree $-3$ or $+1$, the two
cases differing by $\chi_*$ at one end. We fix $c_H$ with $c_H^2=0$ and put
$H=(W_H,c_H)_*$. Then $HB$ has degree $4$, and is not chain homotopic to the
identity if $I(Y_2)\ne0$.
\end{proposition}

\begin{proof}
$H_2(W_H)$ has a basis $R_1,R_2,R_3,F_0$ with $R_i^2=-2$, $R_1\cdot
R_2=R_2\cdot R_3=1$, $F_0\cdot R_2=-1$ and the other products zero. The class
$c=\PD(x_1R_1+x_2R_2+x_3R_3+dF_0)$ is odd on $F_0$ iff $x_2$ is odd, and then
$c^2\equiv2(1+d)\pmod 4$.
\end{proof}

\begin{corollary}\label{cor:grading}
Assume \textup{(H1)}, and let $\F=\F_2$ or $\Q$. If a cobordism $V\colon
Y_{-2}\to Y_2$ with $b_1(V)=b^+(V)=0$ induces an isomorphism on $I(\cdot\,;\F)$,
for instance the mapping cylinder of $\phi$, then $\dim I_j(Y_2;\F)$ does not
depend on $j\in\Z/8$. In particular $\dim I(Y_1;\F)\equiv0\pmod8$ and
$\Delta''_K(1)=0$.
\end{corollary}

\begin{proof}
By (H1), $V_*$ has even degree $-2c^2$ and $H$ odd degree, so $H\circ V_*^{-1}$
is an automorphism of odd degree, and odd numbers generate $\Z/8$. The isomorphism
$I(Y_1)\cong I(Y_2)$ of (H1) has even degree, $\chi(I(Y_1))=\pm2\lambda(Y_1)$
\cite{T1}, and $2\lambda(Y_1)=\Delta''_K(1)$ by Casson's surgery formula.
\end{proof}

\begin{conjecture}\label{conj:B}
Over $\F_2$, $B\simeq g_-\vartheta g_+$, where $g_+\colon C(Y_2)\to C(Y_0,w)$,
$g_-\colon C(Y_0,w)\to C(Y_{-2})$ are the distance-two maps of
\cite[Prop.~3.10]{DLME} for the triads $(Y_2,S^3,Y_0)$, $(Y_0,S^3,Y_{-2})$, both
isomorphisms, and $\vartheta$ is the identification at $Y_0$ fixed by a family
of metrics on $W_B\#(S^2\times S^2)$. So $B\otimes\F_2$ and $B\otimes\Q$ are
isomorphisms.
\end{conjecture}

\step{Status.} \ver: the index counts behind Theorem~\ref{thm:B}, including the
framed index $8E$ of the reducibles on $\nu S$ (close $\nu S$ up in $\CPb$), and
its twisting and degrees; Proposition~\ref{prop:minus4} for simple type, given
\cite[Thms.~4.1, 4.3]{FS2}, and on $K3\#4\CPb$ and $E(4)\#4\CPb$ by computer;
Proposition~\ref{prop:WH}; Corollary~\ref{cor:grading} given (H1), which recovers
the condition $\Delta''_K(1)=0$ of \cite{BL}; the step from $\F_2$ to $\Q$ in
Conjecture~\ref{conj:B} (the cone of $B$ then has finite homology of odd order).
\pla: the chain
property, i.e.\ (H4) (80\%); Proposition~\ref{prop:minus4} in general, by
stretching $L(4,1)$ (75\%). \spe: that the condition modulo $8$ is new. Open:
Conjecture~\ref{conj:B} (85\% that $g_\pm$ are isomorphisms, 60\% that $B\simeq
g_-\vartheta g_+$). Confidence for \S1: 75\%.

\section{Nonvanishing}

\begin{lemma}\label{lem:CH}
Let $A$ be an automorphism of an $r$-dimensional rational vector space, and
$\psi_-$, $\psi_+$ a vector and a covector with $\psi_+(\psi_-)\ne0$. For every
$k\ge1$ the sequence $\psi_+(A^{kM}\psi_-)$, $M\ge0$, has no $r$ consecutive
zeros.
\end{lemma}

\begin{proof}
The characteristic polynomial of $A^k$ has nonzero constant term, so the
sequence satisfies a recurrence of order $r$ that can be solved backwards; $r$
consecutive zeros would propagate to $M=0$.
\end{proof}

Let $g(K)\ge2$. Then $Y_0$ carries a taut foliation with the capped Seifert
surface $\hS$ of genus $g(K)$ as a leaf \cite{Ga}, and $X_0=X^-\cup_{Y_0}X^+$ is the closed
symplectic manifold of \cite[Prop.~15]{KM04}: $H_1(X_0)=H_1(X^\pm)=0$,
$b^+(X^\pm)>0$, and there is $w$ with $\langle w,\hS\rangle=1$. Let
$f_+\colon I(Y_0,w)\to I(Y_2)$ and $f_-\colon I(Y_{-2})\to I(Y_0,w)$ be the maps
of the cobordisms $W_{0\to2}$, $W_{-2\to0}$ of two meridional $2$-handles each,
so that $H=f_+f_-$, and put $C_-=X^-\cup W_{0\to2}\colon\emptyset\to Y_2$ and
$C_+=W_{-2\to0}\cup X^+\colon Y_{-2}\to\emptyset$; they have $H_1=0$ and
$b^+\ge2$ (\ver).

\begin{proposition}\label{prop:caps}
Assume \textup{(H1)} and \textup{(H3)}. There are $z_\pm$ in the algebras
generated by the point class and $H_2(X^\pm)$ with
$\langle\Psi_{C_+}(z_+),H^{-1}\Psi_{C_-}(z_-)\rangle=\pm D^w_{X_0}(z_-z_+)\ne0$ in
$\Q$.
\end{proposition}

\begin{proof}
Gluing along $(Y_0,w)$ gives $\Psi_{C_-}=f_+\Psi_{X^-}$ and
$\Psi_{C_+}=\Psi_{X^+}f_-$, hence the identity. By Mayer--Vietoris, $H_2(X^-)\oplus
H_2(X^+)$ maps onto $\hS^\perp$, and as $H_1(X_0)=0$ and $\hS$ is primitive,
two characteristic classes restrict equally to $\hS^\perp$ iff they differ by
some $2t\PD(\hS)$. By (H3) the restriction of $\mathbf D^w_{X_0}$ to $\hS^\perp$ is
$c\,e^{Q/2}$, $c\ne0$, times a combination of exponentials, one for each coset,
whose coefficients are signed sums of Seiberg--Witten invariants over it. The leaf $\hS$ is symplectic for the form of
\cite[Prop.~15]{KM04}, so the canonical class $\beta_\omega$ has
$\beta_\omega\cdot\hS=2g-2$ and
$(\beta_\omega+2t\PD\hS)^2=\beta_\omega^2+4t(2g-2)$. For $g\ge2$ its
Seiberg--Witten moduli space has negative dimension if $t<0$ and positive
dimension if $t>0$, where simple type kills it. So the coset of $\beta_\omega$
contributes $\pm c\,SW(\beta_\omega)=\pm c\ne0$ \cite{T2}; expand multilinearly.
\end{proof}

\begin{corollary}\label{cor:composites}
Assume \textup{(H1)--(H3)}, $g(K)\ge2$ and that $\phi$ exists, and put
$u=(\phi_*B)^3HB$. For every $p_0$ there are a composite $\Pi_0$ of $\phi_*B$
and $HB$ that begins and ends with at least $p_0$ factors $\phi_*B$ and has at
least three factors $\phi_*B$ between consecutive factors $HB$, and $p'\in8\Z$
with $p'\ge p_0$, such that $\Pi_M=B\,\Pi_0\,u^M(\phi_*B)^{p'}$ satisfies
$\langle\Psi_{C_+}(z_+),\Pi_M\Psi_{C_-}(z_-)\rangle\ne0$ for infinitely many
$M\in8\Z$.
\end{corollary}

\begin{proof}
By (H1)--(H2), $HB$ and $\phi_*B$ are rational automorphisms. Cayley--Hamilton
for $HB$ writes $H^{-1}=B(HB)^{-1}$ as a combination of the $B(HB)^j$, $j<r$, so
by Proposition~\ref{prop:caps} some $B(HB)^j$ pairs nontrivially. As $1$ lies in
the span of the powers $(\phi_*B)^\ell$, $\ell\ge L$, for every $L$, inserting
such combinations between consecutive factors $HB$ and at both ends gives a
composite $B\Pi_0$ of the required form that pairs nontrivially.
Lemma~\ref{lem:CH} for $\phi_*B$ and then for $u$, with $k=8$, gives $p'$ and
infinitely many $M$.
\end{proof}

\step{Status.} \ver: Lemma~\ref{lem:CH}; Proposition~\ref{prop:caps} given
(H1), (H3) and Gabai's leaf (\pla, standard), the topology of $X_0$ and of the
caps having been checked against \cite[Prop.~15]{KM04};
Corollary~\ref{cor:composites} as a deduction. Nothing is known about the parity
of $D^w_{X_0}(z_-z_+)$, so the nonvanishing must be carried over $\Q$, as here;
integrality matters only because Theorem~\ref{thm:family} yields
$2^{n_a-1}\Omega=0$, which is empty modulo two once $n_a\ge2$ (\ver). \pla: (H3)
(85--90\%, conditional on \cite[Hyp.~7.8.1]{FLM}). Open: (H2) (about 65\%). The
genus enters the argument only here, through $\beta_\omega\cdot\hS=2g-2$.
Confidence: 95\% for the deductions, 85\% for Proposition~\ref{prop:caps} with
(H3) as a theorem.

\section{The composition law}

Let $\Pi$ be a composite of $B$, $H$ and maps induced by orientation-preserving
diffeomorphisms between copies of $Y_{\pm2}$, and $X=X(\Pi)$ the closed manifold
obtained by gluing the corresponding copies of $W_B$ and $W_H$ along those
diffeomorphisms and capping with caps $C_\pm$ of the two ends, with $H_1=0$,
$b^+\ge1$ and fixed metrics, carrying classes $z_\pm$ (as in \S2). Let $J_i$, $S_i$ ($1\le i\le n$) be the three-spheres and spheres in
the copies of $W_B$. Let $c_0$ restrict to the classes of the caps, to $c_H$ and to
$0$, and let $E_e$, $e\in\{0,1\}^n$, have $c_1(E_e)=c_0+\sum_ie_i\PD(S_i)$; these
bundles have the same $\kappa$ and pairwise distinct $w_2$ (\ver). Let $g^R_t$,
$t\in Q=[-1,1]^n$, restrict on the $i$-th copy of $W_B$ to the interval $G$ of
\S1 in $t_i$, with the necks along the gluing three-manifolds of length $R$,
except the two inside each copy of $W_B\cup W_H\cup W_B$, which have fixed length
and metric. With $\epsilon(e)=\prod_i(-\epsilon)^{e_i}$, put
\[
\Omega(X;g^R)=\sum_e\epsilon(e)\,\#\bigl\{(t,[A]):t\in\operatorname{int}Q,\
A\text{ a $g^R_t$-instanton on }E_e,\
[A]\in V_{S_1}\cap\dots\cap V_{S_n}\cap\cV(z_-)\cap\cV(z_+)\bigr\}.
\]
Each $J_i$ separates $X$ into two parts with $b^+\ge1$, so all Donaldson and
Seiberg--Witten invariants of $X$ vanish \cite{D90}; at a fixed metric the
problem cut down by $\mu(S_1)\cdots\mu(S_n)$ has virtual dimension $-n$, and the
faces of the cube contribute nothing in total, for the two reasons of \S1.

\begin{theorem}\label{thm:glue}
Assume \textup{(H1)}, \textup{(H4)}, and that consecutive copies of $W_H$ are
separated by at least three copies of $W_B$. For generic auxiliary data in
which the moduli spaces of irreducible instantons on the two halves of each
$W_B$ are transverse to the jumping-line divisors of the capped halves of $S_i$
(Lemma~\ref{lem:nodal}), there is $R_0$ such that for $R\ge R_0$ the count
$\Omega(X;g^R)$ is finite, independent of $R$ and of the data away from the
caps, and $\Omega(X;g^R)=\pm\langle\Psi_{C_+}(z_+),\Pi\,\Psi_{C_-}(z_-)\rangle$. Both sides
vanish unless $c_2(E_0)\in\Z$. If all necks are long, the condition on the
copies of $W_H$ is not needed.
\end{theorem}

\begin{lemma}\label{lem:nodal}
Let $[A_\nu]\in V_{S_i}$ for metrics with $t_i\to-1$ converge to $(A_l,A_r)$,
matched at the trivial connection on $J_i$, with no concentration on $S_i$. Then
$A_l$ or $A_r$ lies in the jumping-line divisor of its capped half of $S_i$.
\end{lemma}

\begin{proof}
On the nodal curve the twist $\cO(-1)$ becomes $\cO(-1)$ on one component and
$\cO$ on the other, and the limiting index-zero operator has kernel
$\{(\sigma_l,\sigma_r)\in H^0(E_l(-1))\oplus H^0(E_r):\sigma_l(p)=g\,\sigma_r(p)\}$,
which vanishes, for every gluing element $g$, iff $E_l\cong E_r\cong\cO\oplus\cO$.
Invertibility passes to large $\nu$ by linear gluing of Cauchy--Riemann
operators.
\end{proof}

\begin{proof}[Sketch of proof of Theorem~\ref{thm:glue}]
Pass to ideal limits as $R\to\infty$ \cite[\S3.9]{KM11}. Index additivity gives
$d=\sum_Pd_P+\sum_\tau\operatorname{ind}\tau+3c$: $d_P$ is the expected dimension
of the cut-down problem on the piece $P$, with its parameters, $\tau$ ranges over
the trajectories on the long necks, and $c$ is the number of matchings at
central flat connections there. A reducible on a copy of $W_B$ lying in $V_S$ is
even on both halves, so it has energy at least $\frac12$ and index $8E-3\ge1$
(Proposition~\ref{prop:energy}); one with $\langle v,S\rangle=0$ meets $V_S$ only
with a point of concentration on $S$; and by Lemma~\ref{lem:nodal} a copy split
along $J_i$ carries the constraint on one half. With the transversality
assumed, $d_P\ge0$ for every piece, except that a copy of $W_B\cup W_H\cup W_B$
with fixed internal necks has $d_P\ge-1$ with one central limit and $d_P\ge-4$
with two. As copies of $W_H$ are three copies of $W_B$ apart, a central neck
borders at most one such piece, and $d\ge2$ once $c>0$. Hence in dimension zero
$c=0$, and the limits are rigid irreducible solutions on the pieces, matched at
irreducible flat connections; each glues uniquely by (H4). The weights factor,
which gives the identity at chain level. Invariance follows from it, as each
piece's map is invariant up to homotopy: $B$ by Theorem~\ref{thm:B}, and $H$
because a reducible on a wall of $W_H$, with its two central matchings, raises
the index by at least $5$.
\end{proof}

\step{Status.} \ver: the index arithmetic, by an exhaustive computation over the
possible limits, and the need for each hypothesis. Without
Lemma~\ref{lem:nodal}, $b$ copies of $W_B$ split along $J$ with trivial halves
reach expected dimension $3-b$; with two copies of $W_B$ between copies of
$W_H$, $k$ adjacent copies of $W_B\cup W_H\cup W_B$ reach $3-k$; without the
transversality, a copy of $W_B\cup W_H\cup W_B$ with two central limits, between
copies of $W_B$ split along $J$ with reducible halves facing it, reaches $0$
(with all necks long, dimension zero needs no transversality). Also \ver: the
algebra of Lemma~\ref{lem:nodal}; the distinct $w_2$. \pla: compactness and
gluing for families of broken metrics whose cuts carry reducibles, excluded in
\cite[\S3.9]{KM11} and handled by index bounds in \cite[\S5]{DLME}; the linear
gluing in Lemma~\ref{lem:nodal}; the transversality, a genericity statement; the
integral sign (85--90\%). Open, and not needed: invariance when the long necks
are shortened. Confidence: 80\% with all necks long, 70\% as stated.

\section{Vanishing in families}

Let $X$ be closed and oriented with $H_1(X;\Z)=0$ and $b^+\ge1$, containing
disjoint separating three-spheres $J_i$ and disjoint spheres $S_i$
($1\le i\le n$) with $S_i\cdot S_i=-4$, $S_i\cap J_i$ a knot and $S_i\cap
J_j=\emptyset$ for $j\ne i$; $[S_i]=F_{l,i}-F_{r,i}$, halves of square $-2$.
Let $(g_t)_{t\in\bar Q}$ stretch $J_i$ as $t_i\to-1$, and $\partial\nu S_i$,
with a metric of positive scalar curvature, as $t_i\to+1$. In the notation of
\cite{FL2b}, let $\ft_e$ be spin$^u$ structures with $w_2(\ft_e)\equiv
w_0+\sum_ie_i\PD[S_i]$, $\Lambda_e=c_1(\ft_e)$ and common $p_1=-4\kappa$, where
$w_0$ is even on the halves and odd on a class orthogonal to the $S_i$, and
$\langle\Lambda_0,S_i\rangle=2$, which is exactly the condition for
$n_a=\frac14(p_1+\Lambda_e^2-\sigma)$ (the $n_D$ of the Strategy) to be
independent of $e$ (\ver). Let $z=\mu(S_1)\cdots\mu(S_n)z'$ with $\deg
z=d_a+n$, $d_a=8\kappa-3(1+b^+)$, and $\Omega=\sum_e\epsilon(e)\#(\bar\cV(z)\cap
\bar M^{w_e}_\kappa(Q))$, as in \S3. Let $\cZ_e=\cV(z)\cap\mathcal
W^{n_a-1}\cap\cM^{*,0}_{\ft_e}(Q)/S^1$ be the cut-down one-manifold of
\cite[(3.62)]{FL2b}, where the sections defining $\mathcal W$, which represent
$\mu_c$ \cite[(3.12)]{FL2b}, sample the spinor on every piece. (H5) consists of:
compactness of $\cZ_e$ over $\bar Q$; transversality for single components, for
the product problem at broken limits (modulo the gauge groups of the pieces and
the diagonal circle, by multivalued perturbations) and for the problem obtained
by forgetting the components that are abelian instantons;
\cite[Prop.~3.29]{FL2b} with parameters; gluing at the lens faces with the sign
$\epsilon$ of Theorem~\ref{thm:B}; and the control of Seiberg--Witten classes on
the positive pieces and the caps used in Theorem~\ref{thm:noreducible}.

\begin{theorem}\label{thm:family}
Assume \textup{(H5)} and $n_a\ge1$, and that the cut-down instanton spaces
defining $\Omega$ are finite sets of regular points over
$\operatorname{int}Q$. Suppose that, for every $e$, no point of the closure of
$\cZ_e$ over $\bar Q$ is reducible on one of the pieces into which the broken
necks cut $X$, other than the $\nu S_i$; that is, the closure meets no stratum
$\iota(M_\fs(\hat\Gamma))\times\Sym^\ell(\hat\Gamma)$ of such a piece
$\hat\Gamma$, at any level $\ell\ge0$, its zero-spinor part included. Then
$\Omega=0$.
\end{theorem}

\begin{proof}[Sketch]
This is the proof of \cite[Thm.~3.33(a)]{FL2b} with faces. By
\cite[Prop.~3.29]{FL2b} the link of the instanton stratum contributes
$\pm2^{n_a-1}\#(\bar\cV(z)\cap\bar M^{w_e}_\kappa(Q))$ ends to $\cZ_e$. A lower
level lowers the dimension by six, and a point of concentration frees
conditions of degree at most four \cite[(3.25)]{FL2b}; a broken $J_i$ lowers it
by four, three for the stabilizer of the trivial connection and one for the
parameter, given the transversality of the product problem. At $t_i=+1$ only
the reducible of energy $\frac14$ on $\nu S_i$ occurs; the exterior problems for
$e$ and $e+e_i$ on $X\setminus\nu S_i$ coincide, since $\PD(S_i)$ vanishes
there, and their ends cancel in the weighted sum. By hypothesis there are no
other ends, and Stokes' theorem gives $2^{n_a-1}\Omega=0$ in $\Q$.
\end{proof}

For $n=0$ this is \cite[Thm.~3.33(a)]{FL2b}, with its hypothesis (3.68), no
reducible at any level, weakened to ``none in the closure of $\cZ$''; at a fixed
metric with spheres it is Theorem~2.1 of the Strategy. Nothing is glued along a
Seiberg--Witten stratum (\ver\ as logic). Reducibles are excluded by energy:

\begin{proposition}\label{prop:energy}
Let the halves be pairwise orthogonal, $\pi$ the orthogonal projection to the
complement of their span, and $v\equiv w_0\pmod2$. Then
$-\frac14v^2\ge-\frac14(\pi v)^2+\frac12\#\{i:\langle v,S_i\rangle\ne0\}$.
\end{proposition}

\begin{proof}
$-v^2=-(\pi v)^2+\frac12\sum_i(\langle v,F_{l,i}\rangle^2+\langle
v,F_{r,i}\rangle^2)$, the evaluations are even as those of $w_0$ are, and one of
them is nonzero if $\langle v,S_i\rangle\ne0$.
\end{proof}

A reducible limit of points of $V_{S_i}$ with $\langle v,S_i\rangle=0$ has a
point of concentration on $S_i$; as the $S_i$ are disjoint, its level satisfies
$\ell\ge T(v)=\#\{i:\langle v,S_i\rangle=0\}$ (\ver).

\begin{theorem}\label{thm:noreducible}
Let $X=X(\Pi_M)$ for the composites of Corollary~\ref{cor:composites}, with the
caps blown up once, $\langle\Lambda_0,e_\pm\rangle$ odd and large, and the family
$g^R$ of \S3, $R\ge R_0(M)$. Assume \textup{(H5)}, with constants independent of
$R$. Then for all large $M\in8\Z$, $n_a\ge1$ and the hypotheses of
Theorem~\ref{thm:family} hold, so $\Omega(X(\Pi_M);g^R)=0$.
\end{theorem}

\begin{proof}[Sketch]
Blowing up does not change $\Psi_{C_\pm}(z_\pm)$ \cite{FS1}, so \S\S2--3 apply,
and Theorem~\ref{thm:glue} gives the finiteness. Let $m$ be the number of copies
of $W_H$. Each copy of $W_B$, with one parameter and one condition of degree two,
raises $8\kappa$ by $1$, and each positive piece $P=X_{-2}(K)\cup W_H\cup(-X_2(K))
\cong\CP^2\#5\CPb$ raises $b^+$, hence $8\kappa$, by $3$. As
$\langle\Lambda,S_i\rangle=2$, $\Lambda$ is even on the halves with difference
two, so the negative pieces $X_{-2}(K)\cup_\phi(-X_2(K))$ add only a bounded
amount to $\Theta=\frac14(\Lambda^2-\sigma)$, while $\Theta(P)=1$. So
$n_a=\Theta-\kappa=\frac18(5m-n)+O(1)$. Proposition~\ref{prop:energy}, and on $P$
the control of Seiberg--Witten classes (a class on the indefinite $P$ may pair
nontrivially with both adjacent spheres at energy $\frac12$ rather than $1$),
give $-\frac14v^2+T(v)\ge\frac12(n-m)-C$; so an unbroken reducible has
$\ell-T(v)\le\frac18(7m-3n)+O(1)$. For $\Pi_M$, $n=4M+O(1)$ and $m=M+O(1)$, so
$n_a\to\infty$ and $\ell-T(v)\to-\infty$. On a chain of reducible components
between broken necks the same inequalities give expected dimension at least
$-2$ (attained by one positive piece between two broken necks); the count of
Theorem~\ref{thm:family} needs at least $-3$, and $-2$ next to a component with
nonzero spinor, granted the transversality after forgetting the abelian
instantons. The caps carry no abelian instanton, as $w$ is odd on $\hS$ and their
periods are generic, and $\langle\Lambda_0,e_\pm\rangle$ large excludes their
Seiberg--Witten classes.
\end{proof}

The same count shows why a family is needed. At a fixed metric
$8\kappa=2n+3m+O(1)$, and $n_a\to\infty$ needs $m/n>\frac25$ while exclusion
needs $m/n<\frac27$; over the cube the range is $\frac15<m/n<\frac37$, and $u$
has $m/n=\frac14$ (\ver). Without $\phi$ only $(HB)^M$ is available, with
$m/n=1$; its pairings are nonzero for infinitely many $M$, by the argument of
Corollary~\ref{cor:composites}, and the exclusion fails, as it must (\ver).

\step{Status.} \ver: Theorem~\ref{thm:family} as a deduction from (H5),
following \cite[\S3]{FL2b}; Proposition~\ref{prop:energy}; the dimension counts
and the ranges above; the energy bound on $P$, a lattice inequality checked by
hand and exhaustively; the bounds on chains, with $-2$ attained;
Theorem~\ref{thm:noreducible} as a deduction. \pla, routine: compactness,
single components, the link with parameters, weighted Stokes. \pla, new: the
transversality of the product problem and after forgetting the abelian
instantons (50--55\%); the lens ends, whose sign must be that of
Theorem~\ref{thm:B} (70\%); the control of Seiberg--Witten classes on $P$
uniformly over $\bar Q$, and the cap estimates (65\%). The caps have $H_1=0$,
which suffices where simple connectivity was assumed, as it enters only through
the goodness of $w_e$ \cite[Def.~3.20]{FL2b} (\pla). Open: a proof of the transversality after forgetting the
abelian instantons (below); gluing at abelian instantons in families, which is
not needed. Confidence that $\Omega(X(\Pi_M))=0$ for large $M$ by this route:
about 30\%.

\section*{The hardest open point}

It is the transversality of the problem obtained by forgetting the abelian
instantons, at broken limits that contain both a component with nonzero spinor
and a chain of abelian instantons. On negative pieces such instantons exist for
every value of the parameters, and their own circle of stabilizers keeps a
cokernel of negative index that no equivariant perturbation removes; so they
can be neither made regular nor excluded by counting, and Feehan and Leness
reach their closed analogue only through gluing hypotheses \cite[Hyps.~7.8.1,
11.3.5]{FLM}. The remaining problem must be made transverse by multivalued
perturbations coupling distant pieces through the phase of the spinor, with a
margin of one dimension. Verdict: the four
statements prove the Main Theorem from (H1)--(H5) (\ver\ as logic; 80\% that no
further compatibility condition is missing). (H1) and (H3) at 85--90\%, (H2) at
about 65\% and (H4) at 80\% are \pla; (H5), at about 30\%, is \spe\ as a whole.
They are positively correlated, and the cosmetic conclusion by this route is
\spe, at about 15\% (10--25\%).

\begin{thebibliography}{DLME}
\bibitem[BL]{BL} S.~Boyer and D.~Lines, \emph{Surgery formulae for Casson's
invariant and extensions to homology lens spaces}, J. Reine Angew. Math.
\textbf{405} (1990), 181--220.
\bibitem[D]{D90} S.~K. Donaldson, \emph{Polynomial invariants for smooth
four-manifolds}, Topology \textbf{29} (1990), 257--315.
\bibitem[DLME]{DLME} A.~Daemi, T.~Lidman and M.~Miller Eismeier,
\emph{Filtered instanton homology and cosmetic surgery},
arXiv:\allowbreak2410.21248.
\bibitem[FL2b]{FL2b} P.~M.~N. Feehan and T.~G. Leness, \emph{PU(2) monopoles.
II}, J. Reine Angew. Math. \textbf{538} (2001), 135--212.
\bibitem[FLM]{FLM} \bysame, \emph{An $\mathrm{SO}(3)$-monopole cobordism
formula relating Donaldson and Seiberg--Witten invariants}, Mem. Amer. Math.
Soc. \textbf{256} (2018), no.~1226.
\bibitem[FS1]{FS1} R.~Fintushel and R.~J. Stern, \emph{The blowup formula for
Donaldson invariants}, Ann. of Math. \textbf{143} (1996), 529--546.
\bibitem[FS2]{FS2} \bysame, \emph{Rational blowdowns of smooth $4$-manifolds},
J. Differential Geom. \textbf{46} (1997), 181--235.
\bibitem[Ga]{Ga} D.~Gabai, \emph{Foliations and the topology of $3$-manifolds.
III}, J. Differential Geom. \textbf{26} (1987), 479--536.
\bibitem[Ha]{Ha} J.~Hanselman, \emph{Heegaard Floer homology and cosmetic
surgeries in $S^3$}, arXiv:\allowbreak1906.06773.
\bibitem[KM04]{KM04} P.~B. Kronheimer and T.~S. Mrowka, \emph{Witten's
conjecture and Property P}, Geom. Topol. \textbf{8} (2004), 295--310.
\bibitem[KM11]{KM11} \bysame, \emph{Khovanov homology is an unknot-detector},
Publ. Math. IH\'ES \textbf{113} (2011), 97--208.
\bibitem[NW]{NW} Y.~Ni and Z.~Wu, \emph{Cosmetic surgeries on knots in $S^3$},
J. Reine Angew. Math. \textbf{706} (2015), 1--17.
\bibitem[T1]{T1} C.~H. Taubes, \emph{Casson's invariant and gauge theory}, J.
Differential Geom. \textbf{31} (1990), 547--599.
\bibitem[T2]{T2} \bysame, \emph{The Seiberg--Witten invariants and symplectic
forms}, Math. Res. Lett. \textbf{1} (1994), 809--822.
\end{thebibliography}

\end{document}
